\originalchapter{耦合振子与正常模}\label{ch:modes}
\originalsection{从交换能量到共同振动}
两个摆用轻弹簧相连, 只拨动其中一个, 不久另一个也振动, 最初的摆却逐渐停下来, 随后过程反转. 相互作用使能量在部件间转移. 若两个摆同向等幅振动, 连结弹簧长度不变; 若反向等幅振动, 连结弹簧不断伸缩. 这两种特殊运动各保持固定的位移比例, 比一般运动容易分析.

\begin{definition}
线性系统的一个正常模保持固定空间形状: $\boldsymbol x(t)=\boldsymbol a q(t)$, 正常坐标满足 $\ddot q+\omega^2q=0$. 非零实频模中, 所有非零分量以同一角频率振动; 零频模允许静止或匀速漂移. 实正常模的分量可以同相或反相, 某些分量可以始终为零.
\end{definition}
两个点质量 $m$ 各用无质量、不可伸长、长度 $\ell$ 的摆线悬挂, 摆球由劲度 $k$ 的轻弹簧相连. 平衡时两摆竖直, 弹簧水平且无预伸长, 只考虑一阶小角运动. $x_1,x_2$ 是相对各自平衡的水平位移. 张力的水平分量给 $-mgx_j/\ell$, 弹簧作用于第一摆为 $k(x_2-x_1)$, 第二摆方向相反. 故
\[
m\ddot x_1=-(mg/\ell+k)x_1+kx_2,\qquad
m\ddot x_2=kx_1-(mg/\ell+k)x_2.
\]
相加、相减引入正常坐标 $q_+=x_1+x_2$, $q_-=x_1-x_2$, 得独立方程
\[
\ddot q_++\omega_+^2q_+=0,\qquad
\ddot q_-+\omega_-^2q_-=0,\qquad
\omega_+^2=g/\ell,\quad\omega_-^2=g/\ell+2k/m.
\]
同向模 $\boldsymbol a_+=(1,1)^{\mathsf T}$ 不拉伸弹簧; 反向模 $\boldsymbol a_-=(1,-1)^{\mathsf T}$ 受到更强回复力, 所以频率较高. 正常坐标把耦合问题变成两种独立简谐运动.

若 $x_1(0)=0$, $x_2(0)=x_0$, 两摆初速均为零, 则
\[
x_1=\frac{x_0}{2}(\cos\omega_+t-\cos\omega_-t),\qquad
x_2=\frac{x_0}{2}(\cos\omega_+t+\cos\omega_-t).
\]
令 $\bar\omega=(\omega_++\omega_-)/2$, $\Delta\omega=\omega_--\omega_+$, 有 $x_1=x_0\sin\bar\omega t\sin(\Delta\omega t/2)$, $x_2=x_0\cos\bar\omega t\cos(\Delta\omega t/2)$. 当耦合弱, $\Delta\omega\ll\bar\omega$, 快振动受到慢包络调制, 称为拍. 包络绝对值的重复周期为 $2\pi/\Delta\omega$, 首次近乎完全转移发生于 $t=\pi/\Delta\omega$. 带符号包络的周期则为 $4\pi/\Delta\omega$. 单摆能量在耦合存在时还涉及弹簧能的分配; “全部能量在某一摆”通常是弱耦合近似, 而不是每一时刻的严格等式.

\originalsection{质量矩阵与一般模态分解}
对 $N$ 个小位移组成的列向量 $\boldsymbol x$, 动能、势能的二次形式为
\[
T=\tfrac12\dot{\boldsymbol x}^{\mathsf T}M\dot{\boldsymbol x},\qquad
V=\tfrac12\boldsymbol x^{\mathsf T}K\boldsymbol x,
\qquad M\ddot{\boldsymbol x}+K\boldsymbol x=0.
\]
$M$ 为对称正定质量矩阵, $K$ 为对称刚度矩阵. 稳定平衡的二次近似要求 $K$ 半正定; 负本征值意味着指数增长而非稳定振动. 零频自由模可以随初速度线性漂移; 要求所有小位移有界时须约束这些零模, 或要求 $K$ 正定. 代入 $\boldsymbol x=\boldsymbol a\cos\omega t$ 得广义本征问题
\begin{equation}\label{eq:generalized}
K\boldsymbol a=\omega^2M\boldsymbol a,\qquad
\det(K-\omega^2M)=0.
\end{equation}
行列式方程给频率, 原线性方程给分量比例. 不应把列向量的每个分量当作可以独立选择的幅度.

为什么正常模足够描述全部运动? 因 $M>0$, 可取实对称平方根 $M^{1/2}$. 令 $\boldsymbol y=M^{1/2}\boldsymbol x$, 方程化为 $\ddot{\boldsymbol y}+D\boldsymbol y=0$, $D=M^{-1/2}KM^{-1/2}$ 为实对称矩阵. 调用有限维实对称矩阵的正交谱定理, 其本征向量构成一组完整正交基. 转回原坐标后可把模态归一为 $\boldsymbol a_r^{\mathsf T}M\boldsymbol a_s=\delta_{rs}$. 展开 $\boldsymbol x=\sum_r\boldsymbol a_r q_r$ 并左乘 $\boldsymbol a_s^{\mathsf T}$, 得 $\ddot q_s+\omega_s^2q_s=0$. 初值的投影为 $q_s(0)=\boldsymbol a_s^{\mathsf T}M\boldsymbol x(0)$ 和 $\dot q_s(0)=\boldsymbol a_s^{\mathsf T}M\dot{\boldsymbol x}(0)$. 零频模应写 $q_s=c_s+v_st$, 而不是用含 $1/\omega_s$ 的公式.

两个不同频率模的质量正交也可直接验证. 对 $K\boldsymbol a_r=\omega_r^2M\boldsymbol a_r$, 与另一模式相减得到 $(\omega_r^2-\omega_s^2)\boldsymbol a_s^{\mathsf T}M\boldsymbol a_r=0$. 简并频率的模可以在该子空间内作质量正交化. $M^{-1}K$ 本身一般不对称, 不能使用普通欧氏内积断言其本征向量正交.

\begin{example}
质量 $2m$ 的物块分别用两根劲度 $k$ 的弹簧与两个质量 $m$ 的物块连接; 整个系统不与固定墙相连, 运动沿弹簧轴. 求全部正常模.
\end{example}
\begin{solution}
以中央 $2m$ 物块为分量1, 另两块为2、3, 有 $M=m\operatorname{diag}(2,1,1)$ 和
\[
K=k\begin{pmatrix}2&-1&-1\\-1&1&0\\-1&0&1\end{pmatrix}.
\]
直接作用于 $(1,1,1)^{\mathsf T}$, $(0,1,-1)^{\mathsf T}$, $(1,-1,-1)^{\mathsf T}$, 可得 $\omega^2=0,k/m,2k/m$. 第一模是整体平移, 不伸缩任何弹簧; 其运动为 $c+vt$. 后两个模分别表示旁边两块反向运动以及中央物块与旁块反向运动. 例如 $K(1,-1,-1)^{\mathsf T}=(4k,-2k,-2k)^{\mathsf T}=(2k/m)M(1,-1,-1)^{\mathsf T}$; 与平移模的质量内积为 $2m-m-m=0$. 一般解实际为
\[
\boldsymbol x=(c+vt)\begin{pmatrix}1\\1\\1\end{pmatrix}
+A\cos(\sqrt{k/m}\,t+\phi)\begin{pmatrix}0\\1\\-1\end{pmatrix}
+B\cos(\sqrt{2k/m}\,t+\psi)\begin{pmatrix}1\\-1\\-1\end{pmatrix}.
\]
三组独立空间向量提供三组位移初值与三组速度初值共六个常数.
\end{solution}

\originalsection{驱动耦合系统与反共振}
若只驱动双摆中的第一个, 外力向量为 $(F_0,0)^{\mathsf T}\cos\omega_dt$. 令 $\alpha=g/\ell+k/m$, $\kappa=k/m$, 无阻尼且避开固有频率时,
\[
\begin{pmatrix}\alpha-\omega_d^2&-\kappa\\-\kappa&\alpha-\omega_d^2\end{pmatrix}
\begin{pmatrix}B_1\\B_2\end{pmatrix}=\begin{pmatrix}F_0/m\\0\end{pmatrix}.
\]
由二元消元或克拉默法则得
\[
B_1=\frac{(F_0/m)(\alpha-\omega_d^2)}{(\omega_+^2-\omega_d^2)(\omega_-^2-\omega_d^2)},\qquad
B_2=\frac{(F_0/m)\kappa}{(\omega_+^2-\omega_d^2)(\omega_-^2-\omega_d^2)}.
\]
接近某一模态频率时, 振幅比接近该模态形状; 在 $\omega_d^2=\alpha$ 时, $B_1=0$ 而 $B_2=-F_0/k$. 外力恰被连结弹簧的力抵消, 被驱动物体的稳态位移反而为零, 称为反共振. 例如零位移、零速度初值时, 特解须补齐次运动:
\[
\begin{split}
x_1(t)&=B_1\cos\omega_dt-\frac{B_1+B_2}{2}\cos\omega_+t
-\frac{B_1-B_2}{2}\cos\omega_-t,\\
x_2(t)&=B_2\cos\omega_dt-\frac{B_1+B_2}{2}\cos\omega_+t
+\frac{B_1-B_2}{2}\cos\omega_-t.
\end{split}
\]
一般初值先取 $\boldsymbol x(0)-\boldsymbol B$ 与 $\dot{\boldsymbol x}(0)$, 按和差投影到两个自由模. 无阻尼时齐次部分不会消失.

若阻尼矩阵为 $\Gamma M$, 每个正常坐标的响应分母为 $\omega_r^2-\omega_d^2+\ii\Gamma\omega_d$. 在实验室分量中, 令 $d_\omega=\alpha-\omega_d^2+\ii\Gamma\omega_d$, 则
\[
\widetilde B_1=\frac{(F_0/m)d_\omega}{d_\omega^2-\kappa^2},\qquad
\widetilde B_2=\frac{(F_0/m)\kappa}{d_\omega^2-\kappa^2}.
\]
有阻尼时第一个分子的虚部通常不为零, 原无阻尼反共振的严格零幅变成凹陷. 一般阻尼矩阵未必与模态同时对角化; 此时必须解完整复矩阵方程, 不能假设所有模独立衰减.
